\clearpage\begingroup
\SourceLeaf{140}\NoSourceHead
\makeatletter\tagsleft@true\let\veqno\@@leqno\makeatother
\newcommand{\Vroot}{\mathord{\vcenter{\hbox{\Large$\surd$}}}}
\newcommand{\Vpoint}{\mathpunct{\vcenter{\hbox{.}}}}
\begin{center}
[V.]\par\bigskip
\textbf{ACHROMATISCHE DOPPELOBJECTIVE, OHNE RÜCKSICHT\\
AUF DICKE UND ABSTAND.}\par\medskip
\rule{.44\linewidth}{.4pt}\par
[Aus Handbuch 21, Bg, S. 29, 30.]\par
\rule{.44\linewidth}{.4pt}
\end{center}
\begin{center}
\begingroup\fontsize{9}{12}\selectfont\setlength{\tabcolsep}{2pt}
\begin{tabular}{lcccccccl}
$(1+m):1$&Brechungsverhältniss&der&rothen&Strahlen&in&der&ersten&Linse;\\
$(1+m'):1$&»&»&»&»&»&»&zweiten&»\quad;\\
$(1+M):1$&»&»&violetten&»&»&»&ersten&»\quad;\\
$(1+M'):1$&»&»&»&»&»&»&zweiten&»\quad;
\end{tabular}\endgroup
\end{center}
\noindent $\frac1\rho,\frac1\sigma,\frac1{\rho'},\frac1{\sigma'}$ Halbmesser der Flächen, positiv, wenn convex;\par
\noindent $p,P$ erste Vereinigungsweite für rothe und violette Strahlen;\par
\noindent\hspace*{.5em}$\omega$\quad Brennweite.

[Dann ist:]
\begin{align}
\rho+\sigma&=\frac1{mp}\tag*{[1]}\\[4pt]
\rho-\sigma&=\frac{2m+4}{m+3}\mathbin{.}\frac1p+\frac{(m+1)\surd(4m+3)}{m(m+3)}\mathbin{.}\frac1p\surd(l-1)\tag*{[2]}\\[4pt]
\rho'+\sigma'&=\frac1{m'}\left(\frac1\omega-\frac1p\right)\tag*{[3]}\\[4pt]
\rho'-\sigma'&=\frac{2m'+4}{m'+3}\left(\frac1\omega+\frac1p\right)+\frac{(m'+1)\surd(4m'+3)}{m'(m'+3)}\left(\frac1\omega-\frac1p\right)\surd(l'-1);\tag*{[4]}
\end{align}
\noindent [Bedingung, dass die sphärische Aberration, die durch beide Linsen des Doppelobjektivs erzeugt wird, verschwindet:]
\begin{equation}
\begin{aligned}
0={}&\frac{(m+1)(4m+3)}{mm(m+3)}\mathbin{.}\frac l{p^3}
+\frac{(m'+1)(4m'+3)}{m'm'(m'+3)}l'\left(\frac1\omega-\frac1p\right)^3\\
&-\frac{4(m'+1)}{m'+3}\mathbin{.}\frac1{\omega p}\left(\frac1\omega-\frac1p\right)\Vpoint
\end{aligned}\tag*{[5]}
\end{equation}

Man setze
\begin{equation}
a=\frac{m(m+1)(4m+3)}{(m+3)(M-m)^3}\tag*{[6]}
\end{equation}
\clearpage\SourceLeaf{141}\HistoricalHead{136}{L}{PHYSIK.\quad NACHLASS.}
\begin{align}
b&=\frac{m'(m'+1)(4m'+3)}{(m'+3)(M'-m')^3}\tag*{[7]}\\[5pt]
c&=\frac{m'(m'+1)}{m'+3}\mathbin{.}\frac{4m(mM'-m'M)}{(M-m)^2(M'-m')^2}\tag*{[8]}\\[5pt]
e&=\frac{2m(m+2)}{m+3}\tag*{[9]}\\[5pt]
f&=\frac{(m+1)\surd(4m+3)}{m+3}\Vroot\frac{a-b+c}{a}\tag*{[10]}\\[5pt]
e'&=\frac{2(m'+2)}{m'+3}\mathbin{.}\frac{2mM'-m'M-mm}{m-M}\tag*{[11]}\\[5pt]
f'&=\frac{(m'+1)\surd(4m'+3)}{m'+3}\Vroot\frac{a-b+c}{b}\Vpoint\tag*{[12]}
\end{align}
\noindent So wird
\begin{align}
\frac{\rho-\sigma}{\rho+\sigma}&=e+f\operatorname{tang}\psi\tag*{[Ia]}\\[5pt]
\frac{\rho'-\sigma'}{\rho'+\sigma'}&=e'+f'\sec\psi,\tag*{[Ib]}
\end{align}
\noindent wo $\psi$ willkürlich angenommen werden kann, wenn bloss die rothen Strahlen genau vereinigt werden sollen.

Will man bloss die violetten Strahlen vereinigen, so setze man
\begin{align}
\frac{\rho-\sigma}{\rho+\sigma}&=E+F\operatorname{tang}\Psi,\tag*{[IIa]}\\[5pt]
\frac{\rho'-\sigma'}{\rho'+\sigma'}&=E'+F'\sec\Psi,\tag*{[IIb]}
\end{align}
\noindent wo die Formeln für $E,F,E',F'$ aus denen für $e,f,e',f'$ entstehen, wenn bloss alle kleinen Buchstaben mit den grossen vertauscht werden. Bestimmt man $\psi$ und $\Psi$ so, dass allen vier Gleichungen zugleich Genüge geschieht, so werden alle Strahlen vereinigt;
\[
\frac{\rho'+\sigma'}{\rho+\sigma}=-\frac{M-m}{M'-m'}\Vpoint
\]
\begin{center}\rule{.25\linewidth}{.4pt}\end{center}
\noindent\hspace*{\parindent}Noch einfacher wird folgendes Verfahren seyn
\[
\frac{M-m}{M'-m'}=\theta,\qquad p=\omega\left(1-\frac{m'}m\theta\right)
\]
\[
\begin{aligned}
a&=\frac{m(m+1)(4m+3)}{m+3}\\[5pt]
b&=\frac{m'(m'+1)(4m'+3)\theta^3}{m'+3}\\[5pt]
c&=\frac{m'(m'+1)}{m'+3}\mathbin{.}4mm\theta\left(1-\frac{m'\theta}m\right)
\end{aligned}
\]
\clearpage\SourceLeaf{142}\HistoricalHead{137}{R}{ACHROMATISCHE DOPPELOBJECTIVE, OHNE RÜCKSICHT AUF DICKE UND ABSTAND.}
\[
\begin{aligned}
e&=\frac{2m(m+2)}{m+3}\\[5pt]
f&=\frac{(m+1)\surd(4m+3)}{m+3}\Vroot\frac{a-b+c}{a}\\[5pt]
e'&=-\frac{2(m'+2)}{m'+3}\mathbin{.}\frac{2m-m'\theta}{\theta}\\[5pt]
f'&=\frac{(m'+1)\surd(4m'+3)}{m'+3}\Vroot\frac{a-b+c}{b}
\end{aligned}
\]
\[
\begin{aligned}
\frac{\rho-\sigma}{\rho+\sigma}&=\frac{2m(m+2)}{m+3}+\frac{(m+1)\surd(4m+3)}{m+3}\Vroot\frac{a-b+c}{a}\mathbin{.}\operatorname{tg}\psi\\[6pt]
\frac{\rho'+\sigma'}{\rho+\sigma}&=-\theta\\[6pt]
\frac{\rho'-\sigma'}{\rho+\sigma}&=+\frac{2(m'+2)}{m'+3}\mathbin{.}m\left(2-\frac{m'}m\theta\right)\\
&\qquad+\frac{(m'+1)\surd(4m'+3)}{m'+3}\mathbin{.}\Vroot\frac{a-b+c}{b}\mathbin{.}\theta\sec\psi
\end{aligned}
\]
\[
1-\frac{M'\theta}M=\frac mM\left(1-\frac{m'}m\theta\right)\Vpoint
\]
\begin{center}[Numerische Berechnung.]\end{center}
\[
\begin{array}{r@{\qquad}r@{\;}c@{\;}l}
m=0{,}504348&M-m&.\,.\,.\,.\,.&8{,}3350164\\
M=0{,}525976&M'-m'&.\,.\,.\,.\,.&8{,}6011905\\
m'=0{,}581810&\theta&.\,.\,.\,.\,.\,.\,.&9{,}7338259\\
M'=0{,}621730&m'&.\,.\,.\,.\,.\,.\,.&9{,}7647812
\end{array}
\]
\[
\begin{array}{r@{\;}l}
\theta m'=&0{,}31521514\\
m-\theta m'=&0{,}18913286\\
{}[\log(m-\theta m')=]&9{,}2767670
\end{array}
\]
\noindent\hspace*{\parindent}[Hier folgen Berechnungen der Konstanten $a,b,c,$ . . . $A,B,C,$ . . ., die schliesslich folgendes Resultat ergeben:]

Also
\[
\begin{aligned}
\frac{\rho-\sigma}{\rho+\sigma}&=0{,}7208547+0{,}8859269\operatorname{tg}\psi\\
&=0{,}7536083+0{,}8974447\operatorname{tg}\Psi\text{[*)]}\\[5pt]
\frac{\rho'-\sigma'}{\rho+\sigma}&=0{,}9974386+1{,}1353112\sec\psi\\
&=1{,}0353119+1{,}1407444\sec\Psi\text{[*)]}.
\end{aligned}
\]
\begingroup\renewcommand{\thefootnote}{}\footnotetext{[*) In der Handschrift steht hier irrtümlich statt $\Psi$ der Buchstabe $\psi'$, der im Vorhergehenden nicht eingeführt ist.]}\endgroup
\noindent Dieser Gleichung geschieht Genüge, wenn man setzt
\par\HistoricalFoot{XI1.}{18}
\clearpage\SourceLeaf{143}\HistoricalHead{138}{L}{PHYSIK.\quad NACHLASS.}
\[
\begin{aligned}
\frac{\rho-\sigma}{\rho+\sigma}&=+1{,}8334367\\[4pt]
\psi&=321^{0}\,28'\,13''{,}40\\
\Psi&=320^{0}\,16'\,8''{,}79\\[4pt]
\frac{\rho'-\sigma'}{\rho+\sigma}&=+2{,}8200047.
\end{aligned}
\]
\begin{center}\rule{.25\linewidth}{.4pt}\end{center}
\clearpage\endgroup
